\documentclass[conference]{IEEEtran}
\usepackage{cite}
\usepackage{amsmath,amssymb,amsfonts}
\usepackage{booktabs}
\usepackage{array}
\usepackage{url}
\usepackage{xcolor}
\usepackage{tabularx}
\newcommand{\CMCA}{\textsc{CMCA}}
\newcommand{\GNW}{\textsc{GNW}}
\newcommand{\IIT}{\textsc{IIT}}
\newcommand{\ORCH}{\textsc{Orch-OR}}
\begin{document}

\title{Where Does Consciousness Lie?\\
\large A Causal Multiscale Consciousness Atlas Research Proposal}

\author{\IEEEauthorblockN{Anonymous Research Proposal}
\IEEEauthorblockA{\textit{Evidence-to-Design Scientific Workflow}\\
Research Plan for Theory-Discriminating Consciousness Science}}

\maketitle

\begin{abstract}
Consciousness is often treated as though it should have a single anatomical location. That framing hides several distinct targets: phenomenal experience, access to information, report, self-awareness, and global state. It also confuses neural correlates with causal mechanisms. This paper synthesizes leading accounts of consciousness and proposes the Causal Multiscale Consciousness Atlas (CMCA), a preregistered design that compares global neuronal workspace, integrated information, recurrent processing, higher-order/predictive, and Orch-OR hypotheses across report-matched and no-report conditions. CMCA combines multiscale recordings, validated content and state proxies, targeted perturbation, effective-connectivity analysis, and out-of-sample theory prediction. Its primary result is a causal evidence vector that classifies evidence as state, content, report, causal necessity, or molecular bridge rather than assigning consciousness to a single coordinate. The Penrose-Hameroff Orch-OR proposal is retained as a direct but contested competitor: a microtubule-related signal would need an independently measured, time- and scale-linked relationship to conscious state that survives neural-network, thermal, measurement, and decoherence null models. The direct conclusion is that current science supports a brain-dependent, dynamically distributed phenomenon but has not settled its causal location or physical mechanism. CMCA offers an ambitious route to locate the relevant process by intervention and theory comparison while preserving an explicit failure path for every claim. This is a design-only proposal and reports no participant, neural, molecular, quantum, or clinical result.
\end{abstract}

\begin{IEEEkeywords}
consciousness, neural correlates, causal perturbation, no-report paradigm, global neuronal workspace, integrated information, recurrent processing, Orch-OR, multiscale neuroscience.
\end{IEEEkeywords}

\section{Introduction}

Consciousness is the fact that there is something it is like to see, hear, feel, remember, imagine, or otherwise experience the world. The question ``Where does consciousness lie?'' is therefore understandable but scientifically underdetermined. A brain region can be active while a person is conscious without that region being sufficient for a particular experience. A verbal report can accompany an experience while also requiring attention, decision, working memory, language, motor preparation, and compliance. Self-awareness can involve a representation of one's own mental state without being a prerequisite for every conscious experience.

The first task is consequently conceptual. This proposal separates five constructs: phenomenal experience, access, report, self-awareness, and conscious state. Phenomenal experience concerns the subjective content; access concerns flexible use of information; report is an observable response channel; self-awareness is a higher-order capacity; and state is the broad capacity for conscious processing. The same neural measurement may be a state marker without being a content marker, or a report marker without being a mechanism marker.

The second task is causal. Koch \textit{et al.} distinguish neural correlates of consciousness from a complete mechanism \cite{koch2016}. A correlate may be necessary, sufficient, downstream, or jointly caused with experience. The field therefore needs perturbations and matched conditions, not only passive localization. Seth and Bayne's review of consciousness theories emphasizes that the major accounts make different biological commitments but are not yet cleanly adjudicated by a single consensus experiment \cite{seth2022}.

The third task is theoretical. Global neuronal workspace (\GNW) emphasizes widespread availability and long-range broadcasting; integrated information theory (\IIT) emphasizes intrinsic causal integration and differentiation; recurrent-processing accounts emphasize feedback loops; higher-order and predictive accounts emphasize metacognitive representation or precision-weighted hierarchical inference. Penrose and Hameroff's orchestrated objective reduction (\ORCH) proposal adds a molecular-scale quantum hypothesis involving microtubules \cite{hameroff2014}. It is a specific and controversial proposal, not a settled fact.

This paper proposes the \emph{Causal Multiscale Consciousness Atlas} (\CMCA). CMCA does not search for a single bright spot. It asks which process is causally required for a named experience, at which spatial and temporal scale, and which theory predicts the joint pattern in held-out data. It combines no-report and report-matched conditions, multiscale recording, targeted perturbation, effective-connectivity modeling, and an independent stress test of the molecular bridge in \ORCH{}. The result is an evidence vector, not a metaphysical consciousness meter.

The contribution is intentionally ambitious. If the proposed convergence occurs, it can identify an experience-sensitive causal network and substantially narrow the theory space. If a candidate signal fails no-report, perturbational, control, or held-out tests, the design localizes the failure rather than allowing a flexible narrative to survive. This is a research plan only: it authorizes no human intervention, laboratory assay, data acquisition, organoid consciousness claim, or observed result.

\section{Survey: The Evidence and the Theory Landscape}

\subsection{From a location question to a causal substrate}

An anatomical location is not the same as a causal substrate. The visual cortex, thalamus, frontoparietal network, and recurrent loops can each be relevant under a particular task or state while playing different roles. A useful answer to ``where'' must specify at least four coordinates: spatial distribution, temporal window, computational operation, and causal status. ``Posterior'' or ``frontal'' without these qualifiers is not a complete scientific claim.

Koch \textit{et al.} review the neural correlates problem and show why the distinction between a neural correlate and a mechanism matters \cite{koch2016}. A candidate signal must be tested against the possibility that it is produced by report, attention, memory, arousal, or a downstream readout. This is especially important in humans, where direct access to private experience is indirect and where the measurement itself can alter the task.

Michel \textit{et al.} call for a maturing science that uses convergent operational definitions, adversarial collaboration, and explicit theory testing \cite{michel2019}. CMCA follows this principle by requiring every theory to state a timing, topology, and intervention prediction before confirmatory data are viewed. A theory does not win because it can explain an outcome after the fact.

\subsection{Global neuronal workspace}

The \GNW hypothesis proposes that information becomes conscious when it is globally available through recurrent broadcasting to multiple cognitive systems. Mashour \textit{et al.} review the relationship among conscious processing, ignition, long-range connectivity, and workspace dynamics \cite{mashour2020}. Dehaene, Lau, and Kouider describe a related distinction between local unconscious processing and globally accessible conscious processing \cite{dehaene2017}.

The strong version of the \GNW prediction is not merely ``frontal activity occurs.'' It is that a late, long-range, cross-system broadcast should accompany conscious access, with a reduced but meaningful form surviving when explicit report is removed. The hypothesis is challenged if its proposed broadcast is entirely explained by response selection, confidence, or task demand, or if a conscious content can be causally established before the predicted global event.

\subsection{Integrated information and complexity}

\IIT begins from the intrinsic causal organization of a system. Tononi \textit{et al.} link consciousness to integrated information and differentiated causal structure \cite{tononi2016}. The attraction of this view is that it treats experience as a property of a system's causal organization rather than as the output of a particular cognitive readout. Its empirical challenge is equally clear: a convenient complexity statistic cannot simply be renamed formal integrated information.

Casali \textit{et al.} introduced the perturbational complexity index (PCI), combining a perturbation with a complexity measure to distinguish levels of consciousness across conditions such as wakefulness, sleep, and anesthesia \cite{casali2013}. A PCI-like measure is useful as a state-level assay of causal capacity. It does not by itself identify which perceptual content is present, locate the content-generating mechanism, or uniquely validate \IIT. CMCA therefore treats complexity as one evidence axis.

\subsection{Recurrent, higher-order, and predictive accounts}

Recurrent-processing accounts emphasize feedback, often in sensory and posterior systems. Their discriminating claim is that a suitably timed local recurrent process may be sufficient for conscious content before later global report-related activity. Higher-order theories emphasize a representation of a mental state by another suitable representation, while predictive-processing approaches emphasize hierarchical generative models, precision, and embodied inference. These families contain variants, so a future study must select explicit model restrictions rather than comparing slogans.

The main empirical challenge is to distinguish recurrent content processing from unconscious recurrence, higher-order representation from confidence, and predictive inference from generic state changes. A report-matched/no-report design and an intervention that targets the proposed process are necessary. A difference in a passive signal is a candidate correlate; it becomes a mechanism candidate only when the intervention changes an experience-sensitive outcome beyond the report pathway.

\subsection{The Orch-OR proposal and its evidentiary boundary}

Hameroff and Penrose propose that consciousness arises from orchestrated objective reductions associated with quantum processes in neuronal microtubules, coupled to neuronal firing and larger-scale dynamics \cite{hameroff2014}. The proposal is more scientifically useful than an unconstrained claim that consciousness is a feature of the universe because it names a molecular substrate and a putative physical process. It is nevertheless controversial.

Tegmark's decoherence analysis argues that quantum states relevant to brain computation would lose coherence on timescales too short for the proposed role \cite{tegmark2000}. Hagan, Hameroff, and Tuszynski challenge aspects of that analysis by changing biological assumptions and estimating longer timescales under their model \cite{hagan2002}. The disagreement does not establish either position. It defines a test burden: a molecular/quantum observable must have a specified timing and scale, be independently measured, track a conscious-state manipulation, and survive neural-network, thermal, sensor, and decoherence alternatives.

A microtubule vibration, high-frequency oscillation, or quantum-like fit is not automatically a consciousness marker. Nor would a negative molecular result prove that quantum effects never occur in biology. It would reject the specified consciousness-generating bridge under the tested conditions. CMCA preserves this granularity.

\subsection{No-report and perturbational evidence}

Tsuchiya \textit{et al.} explain why no-report paradigms are needed: reporting can recruit attention, working memory, decision, and motor processes that are not identical to experience \cite{tsuchiya2015}. No-report does not mean that a private experience is directly observed. It means that the design reduces explicit response demands and uses an independently validated proxy. Proxy validity is a result to establish, not a convenience to assume.

Boly \textit{et al.} review the debate over anterior and posterior neural correlates and show why location disputes can depend on the definition of correlate, content, state, and report \cite{boly2017}. CMCA takes the dispute seriously but changes its form: posterior, frontal, and thalamocortical candidates compete through perturbational predictions, not through a contest over which passive map looks brightest.

\section{Research Gaps and Questions}

The survey identifies six gaps. First, correlation is not cleanly separated from causal necessity. Second, report and metacognition contaminate many candidate signatures. Third, leading theories lack common preregistered discriminating predictions. Fourth, the relevant spatial and temporal scale remains unsettled. Fifth, \ORCH lacks a jointly validated bridge from molecular dynamics to experience. Sixth, state-level complexity markers do not automatically explain content or mechanism.

CMCA asks:
\begin{enumerate}
\item Which distributed process changes a validated experience-sensitive measure when it is perturbed, while report and motor variables are controlled?
\item Which signatures survive report removal and distinguish conscious content from conscious state?
\item Do posterior recurrence, global broadcast, higher-order/predictive dynamics, or integrated causal structure make the best held-out predictions?
\item Is there an independent microtubule-linked molecular signature with the timing and perturbation relationship required by \ORCH?
\end{enumerate}

The central hypothesis is:
\begin{quote}
\emph{Conscious experience is supported by a dynamically integrated causal process with a nontrivial spatial and temporal extent; its theory-specific pattern can be distinguished from report, attention, motor, arousal, and measurement effects through no-report-aware perturbation and held-out prediction.}
\end{quote}

This hypothesis has direct falsifiers. A candidate signature fails as an experience mechanism if it vanishes under no-report control, changes report without changing a validated experience proxy, is reproduced by an unconscious control, or does not predict held-out conditions. The \ORCH bridge fails under this proposal if its molecular signal has no independent experience-linked timing or is explained by neural-network, thermal, measurement, or decoherence null models.

\section{The Causal Multiscale Consciousness Atlas}

\subsection{Evidence vector rather than a single location}

CMCA produces an evidence vector with five axes:
\begin{enumerate}
\item \textbf{State:} whether the system has a broad capacity for conscious processing.
\item \textbf{Content:} which experience is present under matched sensory input.
\item \textbf{Report:} which neural processes support explicit response, confidence, and decision.
\item \textbf{Causality:} which candidate process is necessary for experience-sensitive behavior or proxy.
\item \textbf{Molecular bridge:} whether a specified molecular/quantum process adds explanatory power beyond neural dynamics.
\end{enumerate}

Figure~\ref{fig:cmca} summarizes the proposed flow. It is a decision architecture, not a report of measurements. The distinction between axes is the main protection against overinterpretation. A state score can be high while content localization is weak; a content decoder can be accurate while causal necessity is absent; and a molecular signal can be real while remaining epiphenomenal.

\begin{figure*}[t]
\centering
\fbox{\begin{minipage}{0.92\textwidth}
\centering
\vspace{3pt}
\textbf{D0 construct/preregistration} $\rightarrow$ \textbf{D1 report/no-report content} $\rightarrow$ \textbf{D2 multiscale recording} $\rightarrow$ \textbf{D3 causal perturbation} $\rightarrow$ \textbf{D4 held-out theory comparison} $\rightarrow$ \textbf{D5 independent \ORCH{} bridge test}\\[4pt]
\small State $|$ content $|$ report $|$ causality $|$ molecular bridge are stored as separate evidence axes.\\[5pt]
\textbf{pass} $\rightarrow$ named causal process and theory classification \qquad \textbf{fail/inconclusive} $\rightarrow$ diagnostic code and redesign\\[-1pt]
\small No axis alone establishes a complete theory of consciousness. The present paper authorizes only design, not human or laboratory execution.
\vspace{3pt}
\end{minipage}}
\caption{Causal Multiscale Consciousness Atlas (CMCA). The stages create an auditable path from construct definitions to causal and molecular adjudication.}
\label{fig:cmca}
\end{figure*}

\subsection{Theory commitments}

Table~\ref{tab:theory} translates theory families into future test commitments. These are hypotheses, not observations. The final parameterization and priors must be written by a qualified interdisciplinary team and frozen before confirmatory data are analyzed.

\begin{table*}[t]
\caption{Theory-specific predictions and failure interpretations}
\label{tab:theory}
\centering
\footnotesize
\begin{tabular}{p{0.19\textwidth} p{0.32\textwidth} p{0.34\textwidth}}
\toprule
\textbf{Theory family} & \textbf{Pre-registered target pattern} & \textbf{Failure or boundary condition}\\
\midrule
Global neuronal workspace & Late long-range broadcast/ignition and cross-system availability, with a report-independent component. & Broadcast is fully explained by response, confidence, attention, or motor demands.\\
Integrated information theory & Perturbational differentiation and causal integration track conscious state with a nontrivial causal topology. & A proxy complexity measure is not specific to the predicted causal organization.\\
Recurrent processing & Appropriately timed posterior/sensory feedback is sufficient for conscious content before later report activity. & The same recurrence occurs in unconscious processing or only predicts report.\\
Higher-order/predictive & Metacognitive or precision-weighted hierarchical dynamics add predictive power beyond sensory recurrence. & The effect collapses into confidence, decision, arousal, or an unconstrained model variant.\\
Orch-OR & An independently measured microtubule-linked molecular/quantum signal has predicted timing, scale, and state coupling. & The signal is absent, non-specific, thermally/network explained, or not linked to experience.\\
\bottomrule
\end{tabular}
\end{table*}

\section{Proposed Study Design and Methods}

\subsection{D0: constructs, preregistration, and model registry}

D0 freezes the construct dictionary, conscious/unconscious condition contrasts, report/no-report proxy, state manipulation, target perturbations, recording modalities, quality-control rules, hold-out partition, missing-data policy, and theory restrictions. It also registers a common null model containing sensory processing, arousal, attention, report, motor, and ordinary network dynamics. The common null is not a claim that consciousness is absent; it is the minimum explanation that does not invoke a theory-specific mechanism.

No confirmatory interpretation begins until each theory commits to a timing, spatial topology, intervention response, and expected direction of effect. An exploratory signal may motivate a new study but cannot silently become a preregistered prediction. Condition labels and theory-model labels are blinded until primary analysis code, exclusions, and thresholds are frozen.

\subsection{D1: report-matched and no-report conditions}

D1 uses a future approved perceptual manipulation with sensory-matched seen/unseen or differentiated/undifferentiated content, presented in both report and no-report versions. The design records content, confidence, attention, motor response, arousal, and report timing separately. A no-report proxy must be independently calibrated against the target experience manipulation and cross-checked against a report condition.

The D1 classification is explicit. A difference present only with a button press is a report marker. A difference that separates waking from an unconscious state but not perceptual content is a state marker. A content-specific signal that survives no-report is a candidate experience marker. None becomes a mechanism marker until D3.

\subsection{D2: multiscale recording}

D2 combines complementary scales without assuming that any one modality is privileged. Electrical or magnetic recording supplies temporal dynamics and perturbational responses. Structural and functional imaging supplies spatial registration and network topology. Clinically available fine-scale recordings may be considered only where independently indicated and ethically approved. Each modality receives a signal-quality, sensor-artifact, registration, and vascular-confound record.

The feature vector includes posterior recurrence, frontoparietal broadcasting, thalamocortical coupling where measurable, higher-order/metacognitive variables, effective connectivity, and a PCI-like perturbational complexity measure. D2 identifies candidate regions, network edges, and time windows for intervention. It does not declare a location from a passive association.

\subsection{D3: targeted causal perturbation}

D3 uses a future safety-approved perturbation, such as neuronavigated noninvasive stimulation, with sham and non-candidate target controls. The exact intervention, timing, intensity, and stopping rules must be selected by qualified investigators; this proposal does not execute or prescribe them. The scientific target is a change in an experience-sensitive proxy beyond changes in report, confidence, motor response, or attention.

For candidate process $k$, define the future causal effect
\begin{equation}
\tau_k=E[\phi\mid\operatorname{do}(u_k=1),c]-E[\phi\mid\operatorname{do}(u_k=0),c],
\label{eq:causal}
\end{equation}
where $\phi$ is the validated experience-sensitive proxy, $u_k$ is the approved perturbation of process $k$, and $c$ fixes content and control conditions. The intervention operator $\operatorname{do}(\cdot)$ denotes a causal manipulation rather than a passive association. A mechanism candidate should produce a target-specific change in $\phi$ that is not reproduced by the sham or control target. If it changes report but not $\phi$, the effect is classified as report machinery.

\subsection{D4: effective dynamics and held-out theory comparison}

CMCA represents future multimodal dynamics with a state-space model,
\begin{equation}
x_{t+1}=F_{\theta}x_t+B_{\theta}u_t+\eta_t,\qquad y_t=A_{\theta}x_t+\epsilon_t,
\label{eq:state}
\end{equation}
where $x_t$ is a latent neural state, $u_t$ is the approved perturbation input, $y_t$ is the observed multimodal signal, $F_{\theta}$ is state evolution, $B_{\theta}$ is perturbation coupling, $A_{\theta}$ maps latent state to measurement, and $\eta_t,\epsilon_t$ are process and measurement noise. The matrices and restrictions are theory-dependent. For example, a workspace model constrains long-range broadcast; a recurrent model constrains feedback timing; an \IIT-inspired model constrains causal integration; and a higher-order/predictive model constrains hierarchical or metacognitive contributions.

The primary model comparison uses held-out predictive performance with a frozen flexibility penalty,
\begin{equation}
S_m=\log p(D_{\mathrm{holdout}}\mid M_m)-\lambda C(M_m),
\label{eq:modelscore}
\end{equation}
where $M_m$ is theory model $m$, $D_{\mathrm{holdout}}$ is data not used to choose the restrictions, $C(M_m)$ penalizes flexibility, and $\lambda$ is fixed before analysis. A theory is not selected by in-sample fit or by post hoc storytelling. It must predict condition, timing, topology, and intervention response better than the common null in data held out by participant, session, condition, or site as appropriate.

\subsection{D5: molecular bridge stress test}

D5 asks whether \ORCH{} adds a testable explanatory layer. The future team first registers one or more independently measurable microtubule-linked molecular observables, their predicted timing and scale, the conscious-state contrast, and the measurement method. The molecular signal must be separated from the neural EEG/MEG/fMRI feature used to define the state. It must be evaluated against ordinary network, synaptic, thermal, sensor, and decoherence null models.

The positive criterion is demanding: the molecular signal changes reproducibly with a conscious-state manipulation, has the predicted timing or cross-scale coupling, survives report and measurement controls, and improves held-out theory prediction beyond ordinary neural dynamics. A molecular signal without experience specificity is classified as a molecular correlate or artifact candidate, not as consciousness-generating evidence. Tegmark's decoherence objection and Hagan \textit{et al.}'s assumption-sensitive counteranalysis remain explicit parts of the null registry \cite{tegmark2000,hagan2002}.

\subsection{Identifiability of state, content, report, and mechanism}

The central analytical difficulty is not a shortage of neural measurements. It is that several latent causes can produce a similar observable. Let $z_s$ denote global state, $z_c$ perceptual content, $z_r$ report demand, and $z_m$ the candidate mechanism. A measurement $y$ may be represented as
\begin{equation}
y=h(z_s,z_c,z_r,z_m,a,v)+\epsilon,
\label{eq:identifiability}
\end{equation}
where $a$ represents attention, $v$ represents motor/decision variables, and $\epsilon$ is measurement noise. If all terms change together, the data cannot identify which term generated a signature. D1 holds sensory content and report demands apart; D2 measures complementary spatial and temporal features; and D3 changes a candidate process while observing the other variables. Equation~\eqref{eq:identifiability} is therefore a design statement: identifiability comes from intervention and condition variation, not from a more elaborate post hoc decoder.

CMCA uses nested claims. A state claim requires separation of broad arousal conditions. A content claim requires matched sensory input and an experience-sensitive proxy. A report claim requires prediction of the response channel after content and state are controlled. A mechanism claim requires a target-specific perturbational effect that survives the report controls and predicts held-out observations. This nesting prevents a valid lower-level result from being silently promoted to a higher-level conclusion.

\begin{table*}[t]
\caption{CMCA claim ladder and minimum evidence before promotion}
\label{tab:claimladder}
\centering
\footnotesize
\begin{tabular}{p{0.16\textwidth} p{0.30\textwidth} p{0.37\textwidth}}
\toprule
\textbf{Claim} & \textbf{Minimum evidence} & \textbf{Promotion is blocked when}\\
\midrule
State marker & Reproducible separation of conscious-state conditions using perturbational and passive measures. & The signal follows sensor quality, arousal, or task engagement without a validated state contrast.\\
Content marker & Content-specific result in sensory-matched report and no-report conditions. & The result is explained by motor response, confidence, attention, or report timing.\\
Causal process & Target perturbation changes an experience-sensitive proxy beyond sham/control and replicates out of sample. & The perturbation changes only report, or the control intervention produces the same effect.\\
Theory mechanism & A frozen theory model predicts timing, topology, condition, and intervention data better on hold-out data. & The theory wins only through in-sample fit, flexible post hoc parameters, or a missing null model.\\
Molecular bridge & Independent molecular/quantum signal is time-linked to experience and beats network, thermal, sensor, and decoherence nulls. & The signal is merely present, not experience-specific, or its timing is selected after inspection.\\
\bottomrule
\end{tabular}
\end{table*}

This ladder also clarifies what ``where'' can mean in a scientifically stable report. The answer may be a distributed causal graph with differentiated roles: a posterior loop may supply content, a frontoparietal process may support access, and a thalamocortical or whole-system dynamic may regulate state. Such a result would not be a failure to find a single location. It would be a more precise answer in which location is indexed by the operation and time window under study.

\subsection{Predicted discrimination patterns}

The most informative future result is a dissociation, not a universal increase. For example, if posterior recurrence changes a no-report content proxy while a later frontal signal changes only explicit confidence, recurrent and higher-order interpretations receive different support. If a late broadcast remains necessary after report is removed and predicts flexible cross-system use, the \GNW case strengthens. If a perturbation produces high causal complexity without content selectivity, it supports a state capacity result but not a complete content theory. If several axes are jointly required, a network account becomes more plausible than a single-site account.

The same logic applies to negative results. A failure of no-report proxy validation is a construct failure. A successful proxy with no target-specific perturbation is a correlation-to-causality failure. A causal effect with no held-out theory advantage is a model-identifiability failure. A neural theory pass with an \ORCH{} bridge failure is not a failure of consciousness science; it is evidence against the specified molecular quantum mechanism. Reporting these codes makes the atlas cumulative across laboratories and prevents disagreements from being hidden inside incompatible definitions.

\section{Analysis, Outcome Branches, and Governance}

\subsection{Causal evidence vector}

The primary output is $\mathbf{E}=(E_{\mathrm{state}},E_{\mathrm{content}},E_{\mathrm{report}},E_{\mathrm{causal}},E_{\mathrm{molecular}})$. Each component retains its estimator, control comparison, uncertainty, and interpretation code. The vector is preferable to a single consciousness score because it prevents high state complexity from compensating for a failed content or causal test.

The future team may use a conjunctive mechanism gate,
\begin{equation}
R_{\mathrm{mechanism}}=G_{\mathrm{proxy}}G_{\mathrm{no\mbox{-}report}}G_{\mathrm{causal}}G_{\mathrm{control}}G_{\mathrm{holdout}},
\label{eq:mechanismgate}
\end{equation}
where each indicator is one only when proxy validity, report independence, target-specific causal effect, control rejection, and held-out reproducibility meet their preregistered criteria. A zero identifies a failed link; it does not prove that consciousness is absent.

\begin{table*}[t]
\caption{Future outcome branches and permitted conclusions}
\label{tab:branches}
\centering
\footnotesize
\begin{tabular}{p{0.26\textwidth} p{0.34\textwidth} p{0.27\textwidth}}
\toprule
\textbf{Future pattern} & \textbf{Permitted interpretation} & \textbf{Next action}\\
\midrule
No-report proxy fails validation & The study cannot identify experience independently of report. & Revise construct/proxy; do not localize consciousness.\\
State marker passes, content marker fails & A broad capacity marker is present without content-mechanism evidence. & Separate state from content and test alternative conditions.\\
Content marker passes, causal gate fails & A correlate/decoder exists but necessity is not established. & Audit perturbation, timing, and alternative causal paths.\\
Posterior and frontal candidates both pass different axes & Consciousness may be a distributed process with differentiated roles. & Compare network-level theory predictions rather than choosing a single site.\\
Neural causal gate passes, \ORCH{} bridge fails & A neural mechanism candidate is supported without the molecular quantum bridge. & Continue network theory testing; reject the specified \ORCH{} link under the tested conditions.\\
All neural axes and held-out theory gate pass & A named distributed causal process is supported for the studied experience and conditions. & Replicate across participants/states and test broader content.\\
\bottomrule
\end{tabular}
\end{table*}

\subsection{Reproducibility and data governance}

Any future human study requires an approved consent/data basis, privacy safeguards, ethics review, safety monitoring, preregistered analysis, and qualified investigators. Future data must carry participant/session/condition identifiers, acquisition and preprocessing versions, calibration records, exclusion reasons, perturbation logs, and the unblinding record. At least one participant group, session, modality, or condition should remain held out until theory restrictions and thresholds are fixed. Exploratory and confirmatory analyses remain separate.

The molecular module requires additional governance because claims about quantum biology can be especially vulnerable to sensor coupling, post hoc frequency selection, and theory-flexible interpretation. Independent assay validation, blinded sample/condition labels, temperature and preparation controls, and a predeclared decoherence analysis are necessary. No result from an in-vitro molecular preparation is automatically evidence that a whole organism is conscious.

\subsection{Risks and human review}

The first risk is a report effect masquerading as experience. CMCA counters it with no-report, report-matched, motor, confidence, and attention controls. The second is a state marker masquerading as content or mechanism. The vector and staged gates keep those claims separate. The third is flexible theory fitting. Preregistration, complexity penalties, held-out prediction, and a common null address it. The fourth is overclaiming \ORCH{} from a molecular oscillation. Independent measurement and decoherence/network nulls make the burden explicit.

Qualified neuroscientists must review the perceptual proxy and perturbation target. Qualified clinicians and ethics boards must review any sleep, anesthesia, patient, or stimulation component. Statisticians must review hierarchical and missing-data assumptions. Quantum-information and biophysics experts must review the \ORCH{} observable and decoherence model. The present paper performs none of these interventions and supplies no clinical diagnosis or treatment.

\section{Conclusion}

Current evidence does not place consciousness at one confirmed point in the brain. It supports a stronger and more useful position: conscious experience is brain-dependent and likely involves dynamically integrated processes distributed across spatial and temporal scales, but the causal organization and physical mechanism remain unsettled. Global broadcast, posterior recurrence, higher-order/predictive dynamics, integrated causal structure, and the Orch-OR molecular proposal make different commitments; the field needs experiments that force those commitments into contact with no-report and perturbational evidence.

The Causal Multiscale Consciousness Atlas provides that route. It separates state, content, report, causality, and molecular bridge; tests candidate processes by intervention; evaluates theory models on held-out data; and gives Orch-OR a direct but falsifiable stress test. A future convergence across these axes would justify a precise causal-network claim and could guide neuroscience, brain-inspired computing, and disease research. A failure at any axis would still localize the problem and prevent a correlate, report, or molecular artifact from being promoted to a theory of consciousness. This is an ambitious design, not a conservative disclaimer: it specifies what evidence would earn a stronger answer to where consciousness lies.

\appendices
\section{Idea Evolution and Direction Selection}

The Idea Agent began with passive fMRI localization, asking which region showed the largest difference between conscious and unconscious trials. That seed was rejected as a primary direction because vascular correlation cannot establish causal necessity and because report and task demands contaminate the contrast. A second candidate, a no-report posterior-hot-zone comparison, directly addressed the report confound but still required intervention and theory-specific model comparison. A third candidate, a PCI-like state monitor, supplied a causal complexity measure but could not alone identify content or distinguish \IIT from every other complexity account. A fourth candidate, an Orch-OR molecular bridge stress test, directly engaged the prompt's quantum proposal but risked treating a molecular proxy as whole-organism experience.

CMCA was selected because it combines the first three strengths and retains the fourth as an independent adjudication module. Its selection is traceable to the Survey gap ledger: it turns correlates into interventions, separates report from experience, preregisters theory predictions, estimates spatial and temporal scale, tests the molecular bridge, and preserves the difference between state-level complexity and a mechanism of content.

\section{Evidence Coverage and Required Review}

The Author Agent used only the citation keys registered in the Survey evidence plan. The references support the neural-correlate distinction, theory landscape, global workspace, integrated information, no-report methods, perturbational complexity, anterior/posterior debate, Orch-OR, and decoherence disagreement. The report does not claim that any theory has won, that a molecular quantum process has been observed, or that a participant study has been completed.

The requested dual-engine search was attempted through OpenAlex and AnySearch with three focused queries. Both connections were refused by the local network (`WinError 10061`), so this workflow does not claim an exhaustive live literature census. The Seth-Bayne Nature Reviews Neuroscience page was verified through the specified in-app browser for title, authors, date, journal, volume, pages, and abstract. Full-text applicability, current experimental status, and all platform-specific methodological assumptions remain human-review items before real execution.

\begin{thebibliography}{99}

\bibitem{koch2016}
C. Koch, M. Massimini, M. Boly, and G. Tononi, ``Neural correlates of consciousness: progress and problems,'' \emph{Nature Reviews Neuroscience}, vol. 17, pp. 307--321, 2016, doi: 10.1038/nrn.2016.22.

\bibitem{seth2022}
A. K. Seth and T. Bayne, ``Theories of consciousness,'' \emph{Nature Reviews Neuroscience}, vol. 23, pp. 439--452, 2022, doi: 10.1038/s41583-022-00587-4.

\bibitem{mashour2020}
S. L. Mashour \textit{et al.}, ``Conscious processing and the global neuronal workspace hypothesis,'' \emph{Neuron}, vol. 105, pp. 776--798, 2020, doi: 10.1016/j.neuron.2020.01.026.

\bibitem{dehaene2017}
S. Dehaene, H. Lau, and S. Kouider, ``What is consciousness, and could machines have it?'' \emph{Science}, vol. 358, pp. 1211--1214, 2017, doi: 10.1126/science.aan8871.

\bibitem{tononi2016}
G. Tononi \textit{et al.}, ``Integrated information theory: from consciousness to its physical substrate,'' \emph{Nature Reviews Neuroscience}, vol. 17, pp. 450--461, 2016, doi: 10.1038/nrn.2016.44.

\bibitem{casali2013}
M. M. Casali \textit{et al.}, ``A theoretically based index of consciousness independent of sensory processing and behavior,'' \emph{Science Translational Medicine}, vol. 5, 198ra105, 2013, doi: 10.1126/scitranslmed.3006294.

\bibitem{tsuchiya2015}
N. Tsuchiya \textit{et al.}, ``No-report paradigms: extracting the true neural correlates of consciousness,'' \emph{Trends in Cognitive Sciences}, vol. 19, pp. 757--770, 2015, doi: 10.1016/j.tics.2015.10.002.

\bibitem{boly2017}
M. Boly \textit{et al.}, ``Are the neural correlates of consciousness in the front or in the back of the cerebral cortex?'' \emph{Journal of Neuroscience}, vol. 37, pp. 9603--9613, 2017, doi: 10.1523/JNEUROSCI.3218-16.2017.

\bibitem{michel2019}
M. Michel \textit{et al.}, ``Opportunities and challenges for a maturing science of consciousness,'' \emph{Nature Human Behaviour}, vol. 3, pp. 104--107, 2019, doi: 10.1038/s41562-019-0531-4.

\bibitem{hameroff2014}
S. Hameroff and R. Penrose, ``Consciousness in the universe: a review of the Orch OR theory,'' \emph{Physics of Life Reviews}, vol. 11, pp. 39--78, 2014, doi: 10.1016/j.plrev.2013.08.002.

\bibitem{tegmark2000}
M. Tegmark, ``Importance of decoherence in brain processes,'' \emph{Physical Review E}, vol. 61, pp. 4194--4206, 2000, doi: 10.1103/PhysRevE.61.4194.

\bibitem{hagan2002}
S. Hagan, S. Hameroff, and J. A. Tuszynski, ``Quantum computation in brain microtubules: decoherence and biological feasibility,'' \emph{Physical Review E}, vol. 65, 061901, 2002, doi: 10.1103/PhysRevE.65.061901.

\end{thebibliography}

\end{document}
